\chapter{Transformations}
\vspace{-1cm}
Up to this point, we've just been drawing functions and finding key details about them (asymptotes, domain, range). Now it's time to move functions around and warp them in all kinds of ways. You probably learned how to do this back in Algebra II (or Integrated III). The goal here is to really cement that understand.

\noindent If $f(x)$ is the parent function, these are what the transformations look like:
\[
    af\left(\frac{1}{b}(x - h)\right) + k
\]
\begin{itemize}
    \item $h$ - Horizontal shift (positive = right)
    \item $k$ - Vertical Shift (positive = up)
    \item $a$ - Stretching and Shrinking along the vertical axis
    \begin{itemize}
        \item Larger Values = More Vertical Stretching
    \end{itemize}  
    \item $b$ - Stretching and Shrinking along the horizontal axis
    \begin{itemize}
        \item Larger Values = More Horizontal Stretching
    \end{itemize}
\end{itemize}
\textit{Note: Some texts use $b$ instead of $\dfrac{1}{b}$. When using $b$, larger values of $b$ horizontally \textbf{compress} instead of stretch.}

Let's try to understand why each of these variables transform the function the way they do.